\podnaslov{Opisna statistika}

Tukaj \pojem{statistika} pomeni povzetek opaženih vrednosti; če opazimo $x_1,
\ldots, x_n$, potem je statistika $f(x_1, \ldots, x_n)$.
Številu $n$ pravimo \pojem{numerus}.
Za empirično porazdelitev definiramo \pojem{aritmetično sredino} kot povprečje
vrednosti $x_i$, \pojem{empirični standardni odklon} pa kot
\[
  \sigma = \sqrt{\inv{n} \sum_{k=1}^n (x_k - \overline{x})^2 }.
\]
Če uredimo vrednosti kot $x_{(1)} \le x_{(2)} \le \cdots \le x_{(n)}$, pravimo,
da je $x_{(k)}$ \pojem{$k$-ta vrstilna statistika}.
Za neko verjetnost $\alpha \in (0,1)$ je \pojem{kvantil} slučajne spremenljivke
$X$ tako število $x_\alpha$, da je $P(X < x_\alpha) \le \alpha \le P(X \le
x_\alpha)$.
Kvantil ni enolično določen, razen v primeru, ko je gostota porazdelitve strogo
pozitivna na nekem intervalu, drugje pa enaka $0$.
Nekaj kvantilov ima tudi svoja imena:
\begin{itemize}
\item $x_{1/2}$ je \pojem{mediana},
\item $x_{1/3}$ in $x_{2/3}$ sta prvi in drugi \pojem{tercil},
\item $x_{1/4}, x_{1/2}, x_{3/4}$ so \pojem{tercili}
\item desetine dajo \pojem{decile}, stotine \pojem{centile}, itd.
\end{itemize}
Če kvantili niso enolično določeni, jim določamo arbitrarne vrednosti.
Mediani običajno pripišemo vrednost
\[
  \pol (x_{1/2}^{\text{min}} + x_{1/2}^{\text{max}}).
\]
Povezan pojem je \pojem{interkvartilni razmik} (interquartile range)
\[
  \text{IQR} = x_{3/4}^{\text{max}} - x_{1/4}^{\text{min}}.
\]

\vprasanje{Kaj so kvantili? Kaj je IQR?}

Številu pojavitev posamične vrednosti pravimo \pojem{frekvenca}.
Frekvence lahko sestavimo skupaj in prikažemo v histogramu.
Želeli bi si, da le-ta čim manj odstopa od teoretične gostote porazdelitve.
V primeru Gaussove gostote po kriteriju najmanjših kvadratov pridemo do
Scottovega pravila
\[
  \text{širina} \approx \sigma \sqrt[3]{\frac{24 \sqrt{2 \pi}}{n}},
\]
kar pa veselo ignoriramo, ko je to pripravno.

\vprasanje{Kako postaviš dober histogram?}

Po dogovoru so \pojem{osamelci} vrednosti izven intervala $[x_{1/4}^{\text{min}}
- \frac{3}{2} \text{IQR}, x_{3/4}^{\text{max}} + \frac{3}{2} \text{IQR}]$.


% LocalWords:  kvantilov tercil tercili decile centile kvantili kvartili IQR
% LocalWords:  interkvartilni
